\begin{enumerate}
\item $A = \dfrac 12 - \dfrac 12 \times \dfrac 43 = \dfrac 12 - \dfrac23 = \dfrac{1 \times 3 - 2 \times 2}{2 \times 3} = \dfrac{-1}{6}$.
\item 2 croissants coûtent 3~euros donc $2 \times 5 = 10$ croissants coûtent $3  \times 5 = 15$~euros.
\item Multiplier un prix par deux, c'est le multiplier par $1 + 1 = 1 + \dfrac{100}{100}$, c'est donc l'augmenter de $100~\,\%$.
\item Augmenter un prix de 10\,\%, c'est le multiplier par $1 + \dfrac{10}{100} = 1 + 0,10 = 1,1$.

Le prix initial était donc de $\dfrac{110}{1,10} = 100$. l'augmentation est donc de 10~\euro.
\item La masse de \np{1000} millilitres d'huile est égale à 900~grammes, donc la masse de 250  millilitres d'huile est égale à 225~grammes (table de 4), et enfin la masse de $3 \times 250$ millilitres d'huile est égale à $3 \times 225 = 675$~grammes, ou 0,675~kg.
\item Le coefficient directeur de la droite (AB) est égal à $\dfrac{y_{\text{B}} - y_{\text{A}}}{x_{\text{B}} - x_{\text{A}}} = \dfrac{106 - 100 }{4 - 1} = \dfrac 63 = 2$.
\item Le coefficient directeur de la droite (AB) est égal à $\dfrac{y_{\text{B}} - 4}{1 -  0} = \dfrac{y_{\text{B}} - 4}{1} = y_{\text{B}} - 4 = - 0,1$, d'où :
$y_{\text{B}} = 4 - 0,1 = 3,9$.
\item $(x - 3)(x + 2) = x^2  + 2x - 3x - 6 = x^2 - x - 6$.
\item $V = \dfrac 13 \pi r^2 h$, donc $3V = \pi r^2 h$ et $h = \dfrac{3V}{\pi r^2}$.
\item On a $f(1) = 2 - 5 + 3 = 5 - 5 = 0$, donc 1 est un antécédent de 0.
\item Avec des notations évidentes, on a :

$m_{\text{A}} = \dfrac{9 + 10 + 10 + 11}{4} = \dfrac{40}{4} = 10$.

$m_{\text{B}} = \dfrac{7 + 10 + 10 + 13}{4} = \dfrac{40}{4} = 10$.

$\sigma_{\text{A}}^2 = \dfrac{(- 1)^2 + 0 + 0 + 1^2}{4} = \dfrac{2}{4} = \dfrac12$.

$\sigma_{\text{B}}^2 = \dfrac{(- 3)^2  + 0 + 0 + 3^2}{4} = \dfrac{18}{4} = 4,5$. Réponse D.
\end{enumerate}

\medskip
